\subsection{Inverse of a Linear Function} 
\begin{theorem}
Every linear function is one-to-one, and its inverse is also a linear function.
\end{theorem}
This is easiest to see because of the \makeblank, and because the graph of the inverse is a \makeblank[1in] of the graph of the function.
\begin{example}
Let $f(x)=2x-8$. Find $f\inv$.
\begin{lpic}{}{8}
\twocolleft{-1}{8}{Step}{Why?}
\regtextleft{1}{-2}{$x=2y-8$};
\regtextleft{1}{-8}{$f\inv(x)={}$};
\end{lpic}
\end{example}
\begin{example}
Let $f(x)=\frac{2}{5}x+7$. Find $f\inv$.
\begin{lpic}{}{8}
\twocolleft{-1}{8}{Step}{Why?}
\regtextleft{1}{-2}{$x=\frac{2}{5}y+7$};
\regtextleft{1}{-8}{$f\inv(x)={}$};
\end{lpic}
\end{example}